\section{Data}
\label{sec:data}

\subsection{The \figsum{} gold standard}
\figsum{} \citep{ramesh2015figsum} provides two gold standards. \textbf{GS1} lists, for each figure, every
sentence of the article that annotators judged to be associated with it; this is our task. \textbf{GS2}
contains sentences selected for a short figure \emph{summary}; we use it only to check our reproduction of the
published baselines and never for tuning.

\paragraph{How GS1 was built.} According to the paper, four biologists annotated two articles each (eight
articles), and two of them, B1 and B2, both annotated eleven further articles, marking 303 and 383 sentences
respectively. The released gold is a single combined list of 678 sentences for 94 figures in 19 articles.

\paragraph{Agreement.} On the doubly annotated articles the two annotators agree at Cohen's $\kappa=0.68$
\citep{cohen1960kappa}. Kappa corrects raw agreement for the agreement expected by chance:
\[
\kappa=\frac{P_{\text{agree}}-P_{\text{chance}}}{1-P_{\text{chance}}},
\]
where $P_{\text{agree}}$ is the fraction of (figure, sentence) decisions on which both annotators gave the same
yes/no answer, and $P_{\text{chance}}$ is the agreement expected if each answered at random with their own
rate of ``yes''. The correction matters because most sentences are not associated with a given figure. As an
illustration (the numbers are ours, not from the paper): if each annotator says ``yes'' for 10\% of sentences,
then $P_{\text{chance}}=0.1\cdot0.1+0.9\cdot0.9=0.82$, and a raw agreement of 94\% gives
$\kappa=(0.94-0.82)/0.18=0.67$. A value of 0.68 is conventionally called substantial agreement
\citep{landis1977}, but it also means that the gold standard is far from unambiguous, and that no system
should be expected to approach an F1 of 1.

\paragraph{Full text.} The gold standard lists sentences but not the articles themselves. We obtained the full
text of the 18 articles that have gold labels, mostly from the PubMed Central open-access collection, parsed
them, and matched the gold sentences to the parsed text. This gives \textbf{18 articles, 91 figures and 657
gold sentences}.

\subsection{The fixed split}
\label{sec:split}
\begin{table}[h]
\centering
\small
\begin{tabular}{lrrrrrrr}
\toprule
Split & Papers & Figures & Sentences & Gold & Explicit & Near & Far\\
\midrule
Train (5 folds of 3 papers) & 15 & 73 & 2{,}232 & 549 & 206 & 231 & 112\\
Test (sealed)               &  3 & 18 &   463 & 108 &  32 &  46 &  30\\
\bottomrule
\end{tabular}
\caption{The split used for all learned models. ``Sentences'' counts body sentences only (abstract to
conclusion), without captions, tables or references; these are short biomedical journal articles of roughly
2{,}500 to 5{,}000 words.}
\label{tab:split}
\end{table}

The three test papers were chosen from the five articles added last, which no earlier experiment had
touched. They turn out to be harder than the training papers: 28\% of their gold is far from any citation,
against 20\% in training. The training folds were balanced greedily by number of figures.

\subsection{Where the implicit references are}
Figure~\ref{fig:dist} shows how far the implicit gold sentences of the training papers lie from the nearest
citation of their figure. Of 343 implicit sentences, 176 (51\%) are exactly one sentence away, 55 are two or
three sentences away, and 112 (33\%) are further. More than half of the far sentences (60 of 112) are in the
Discussion. This distribution explains much of what follows: a rule that looks one sentence either side of a
citation is hard to beat, and the far references are both the hardest and the least represented.

\begin{figure}[h]
\centering
\begin{tikzpicture}
\begin{axis}[
  ybar, bar width=16pt, bar shift=0pt, width=0.82\linewidth, height=5.2cm,
  ymin=0, ymax=200, xmin=0.4, xmax=5.6, ylabel={Implicit gold sentences},
  xlabel={Distance to the nearest citation of the figure (sentences)},
  xtick={1,2,3,4,5}, xticklabels={1,2,3,4--10,{more than 10}},
  nodes near coords, nodes near coords style={font=\small},
  axis lines*=left, ymajorgrids, grid style={soft},
  tick label style={font=\small}, label style={font=\small}]
\addplot[fill=steel,draw=steel] coordinates {(1,176) (2,43) (3,12)};
\addplot[fill=accent,draw=accent] coordinates {(4,20) (5,92)};
\end{axis}
\end{tikzpicture}
\caption{Distance from each implicit gold sentence to the nearest explicit citation of its figure
(15 training papers, 343 implicit sentences). Blue: \emph{near} (231). Orange: \emph{far} (112).}
\label{fig:dist}
\end{figure}

\subsection{Other data}
\begin{itemize}
  \item \textbf{arXiv-28:} 28 arXiv papers with \LaTeX{} source. The \verb|\ref| commands give the true
        explicit references for free. Used to develop and measure the explicit grammar.
  \item \textbf{bioRxiv-177:} 177 bioRxiv papers with both XML and PDF. Used as a held-out benchmark for
        explicit references.
  \item \textbf{PMC-6049:} 6{,}049 open-access PubMed Central papers, used to mine free training labels
        (Section~\ref{sec:fakelabels}).
\end{itemize}
